% Plantilla de presentación - IIIEPE
% Alumno: Mtro. Martin Jonathan de la Cruz Muñoz
% Matrícula: 00874
%
% ==============================================================================
% PLANTILLA MAESTRA DE PRESENTACIONES - IIIEPE / MAEM
% ==============================================================================
% Uso recomendado:
% 1. Duplica este archivo y renombra la copia por materia, tetramestre y actividad.
%    Ejemplo: MAEM_T1_Fundamentos_Actividad01.tex
% 2. Actualiza \title, el título visible de portada y el contenido de las
%    diapositivas de la sección "Plantilla de actividad".
% 3. Conserva las diapositivas de identidad del programa cuando la actividad
%    requiera contextualizar el posgrado; si no son necesarias, puedes comentar
%    esas secciones con %.
% 4. No pegues instrucciones completas de la consigna en el producto final. Usa la
%    consigna como lista de cumplimiento y redacta una respuesta académica propia.
% 5. Cada actividad debe cerrar con evidencia de aprendizaje: conclusión, postura,
%    implicación didáctica o propuesta de mejora.
% 6. Si la actividad pide cuadro, esquema, mapa, línea de tiempo, matriz, tabla,
%    diagrama o ruta de intervención, constrúyelo en LaTeX/TikZ cuando sea viable.
%
% Enfoque académico para actividades de la MAEM:
% - Partir de un problema de enseñanza/aprendizaje de las matemáticas.
% - Relacionar teoría, práctica docente, contexto escolar y evidencia.
% - Evitar resúmenes planos: definir -> explicar función -> analizar límite ->
%   proponer aplicación -> valorar efecto.
% - Integrar tecnología educativa, evaluación formativa, resolución de problemas
%   e innovación solo cuando realmente aporten al objetivo de la actividad.
% - Priorizar una voz profesional: clara, reflexiva, argumentada y situada.
%
% Estructura sugerida por actividad:
% [ ] Portada con materia, actividad, alumno, matrícula, docente y fecha.
% [ ] Objetivo de la actividad redactado en una frase propia.
% [ ] Problema o pregunta detonadora.
% [ ] Desarrollo con conceptos, criterios, evidencias y análisis.
% [ ] Producto solicitado visible: tabla, esquema, diapositiva, matriz, etc.
% [ ] Aplicación a la práctica docente o a la enseñanza de las matemáticas.
% [ ] Conclusión con postura, razón y consecuencia.
% [ ] Fuentes consultadas en formato consistente.
%
% Checklist antes de entregar:
% [ ] El título de la presentación y la actividad están actualizados.
% [ ] La materia y el tetramestre corresponden al curso real.
% [ ] El contenido responde exactamente a la consigna.
% [ ] La malla/identidad MAEM no sustituye el desarrollo de la actividad.
% [ ] Hay análisis propio; no solo copia o paráfrasis superficial.
% [ ] Las diapositivas no están saturadas; dividir contenido si es necesario.
% [ ] Ortografía, acentos, nombres propios y fechas revisadas.
% [ ] El PDF compila sin errores y las imágenes institucionales cargan.
% ==============================================================================

\pdfminorversion=7
\documentclass[
	spanish,
	aspectratio=1610,
	xcolor={dvipsnames,table}
]{beamer}

% Lienzo 16:10 ampliado: da más espacio por diapositiva sin cambiar el diseño.
\geometry{paperwidth=19.2cm,paperheight=12cm}

% Idioma, tipografía y recursos visuales
\usepackage[utf8]{inputenc}
\usepackage[T1]{fontenc}
\usepackage[spanish,es-nodecimaldot]{babel}
\usepackage{lmodern}
\usepackage{booktabs}
\usepackage{graphicx}
\usepackage{ragged2e}
\usepackage{tikz}
\usepackage{hyperref}

% Datos del alumno y del programa
\newcommand{\studentname}{Mtro. Martin Jonathan de la Cruz Muñoz}
\newcommand{\studentshort}{M. de la Cruz Muñoz}
\newcommand{\studentid}{00874}
\newcommand{\studentemail}{martin.delacruz@campus.iiiepe.edu.mx}
\newcommand{\programname}{Maestría en Aprendizaje y Enseñanza de las Matemáticas}
\newcommand{\programshort}{MAEM}
\newcommand{\modality}{Modalidad presencial}
\newcommand{\periodname}{1er tetramestre marzo--julio 2026}

% Datos institucionales
\newcommand{\institutionname}{IIIEPE}
\newcommand{\institutionlongname}{Instituto de Investigación, Innovación y Estudios de Posgrado para la Educación del Estado de Nuevo León}
\newcommand{\institutionurl}{https://iiiepe.edu.mx/}
\newcommand{\institutionplace}{Monterrey, Nuevo León}

% Recursos oficiales descargados desde el sitio del IIIEPE
\newcommand{\iiiepelogo}{img/iiiepe/logo-iiiepe-blanco.png}
\newcommand{\iiiepemark}{img/iiiepe/marca-iiiepe.png}
\newcommand{\iiiepecover}{img/iiiepe/portada-institucional.png}
\newcommand{\iiiepedots}{img/iiiepe/logo-puntos.png}
\newcommand{\maemplanimage}{img/iiiepe/maem-plan-estudios.png}
\newcommand{\maemsource}{https://iiiepe.edu.mx/maestria-en-aprendizaje-y-ensenanza-de-las-matematicas/}

% Paleta inspirada en la identidad visual del sitio institucional
\definecolor{iiiepeNavy}{HTML}{163B5C}
\definecolor{iiiepeBlue}{HTML}{1F6F9F}
\definecolor{iiiepeAqua}{HTML}{20B7A8}
\definecolor{iiiepeOrange}{HTML}{F28C28}
\definecolor{iiiepePaper}{HTML}{F6F8F7}
\definecolor{iiiepeInk}{HTML}{1E2A32}
\definecolor{iiiepeMuted}{HTML}{667085}

\mode<presentation>{
	\usetheme{default}
	\usefonttheme{professionalfonts}
	\setbeamertemplate{navigation symbols}{}
	\setbeamertemplate{blocks}[rounded][shadow=false]
	\setbeamercovered{invisible}
}

\hypersetup{
	colorlinks=true,
	urlcolor=iiiepeBlue,
	linkcolor=iiiepeNavy
}

\setbeamercolor{background canvas}{bg=white}
\setbeamercolor{normal text}{fg=iiiepeInk,bg=white}
\setbeamercolor{structure}{fg=iiiepeBlue}
\setbeamercolor{frametitle}{fg=white,bg=iiiepeNavy}
\setbeamercolor{title}{fg=white,bg=iiiepeNavy}
\setbeamercolor{subtitle}{fg=white}
\setbeamercolor{block title}{fg=white,bg=iiiepeBlue}
\setbeamercolor{block body}{fg=iiiepeInk,bg=iiiepePaper}
\setbeamercolor{alerted text}{fg=iiiepeOrange}
\setbeamerfont{title}{size=\LARGE,series=\bfseries}
\setbeamerfont{frametitle}{size=\large,series=\bfseries}
\setbeamerfont{block title}{series=\bfseries}
\setbeamertemplate{itemize item}{\textcolor{iiiepeAqua}{\large$\blacktriangleright$}}
\setbeamertemplate{itemize subitem}{\textcolor{iiiepeOrange}{\scriptsize$\blacksquare$}}
\setbeamertemplate{enumerate item}{\textcolor{iiiepeBlue}{\insertenumlabel.}}

\setbeamertemplate{frametitle}{
	\nointerlineskip
	\begin{beamercolorbox}[wd=\textwidth,ht=1.05cm,dp=0.22cm,leftskip=0.35cm,rightskip=0.25cm]{frametitle}
		\begin{minipage}[c]{0.72\textwidth}
			\insertframetitle
		\end{minipage}
		\hfill
		\raisebox{-0.1cm}{\includegraphics[height=0.58cm]{\iiiepelogo}}
	\end{beamercolorbox}
	\vspace{-0.04cm}
	{\color{iiiepeOrange}\rule{\textwidth}{1.4pt}}
}

\setbeamertemplate{footline}{
	\leavevmode%
	\hbox{%
		\begin{beamercolorbox}[wd=.34\paperwidth,ht=2.7ex,dp=1ex,leftskip=1em]{author in head/foot}%
			\color{white}\studentshort
		\end{beamercolorbox}%
		\begin{beamercolorbox}[wd=.38\paperwidth,ht=2.7ex,dp=1ex,center]{title in head/foot}%
			\color{white}\programshort\ · \modality
		\end{beamercolorbox}%
		\begin{beamercolorbox}[wd=.20\paperwidth,ht=2.7ex,dp=1ex,rightskip=1em plus 1fil]{date in head/foot}%
			\color{white}Matrícula \studentid\hfill\insertframenumber/\inserttotalframenumber
		\end{beamercolorbox}%
	}%
	\vskip0pt%
}
\setbeamercolor{author in head/foot}{bg=iiiepeNavy}
\setbeamercolor{title in head/foot}{bg=iiiepeBlue}
\setbeamercolor{date in head/foot}{bg=iiiepeNavy}

\newcommand{\accentbar}{%
	\begin{tikzpicture}[remember picture,overlay]
		\path[use as bounding box] (0,0) rectangle (0,0);
		\fill[iiiepeAqua] (current page.south west) rectangle ([xshift=0.42\paperwidth,yshift=0.08cm]current page.south west);
		\fill[iiiepeOrange] ([xshift=0.42\paperwidth]current page.south west) rectangle ([xshift=\paperwidth,yshift=0.08cm]current page.south west);
	\end{tikzpicture}
}

\newcommand{\sectionframe}[1]{%
	\begin{frame}[plain]
		\begin{tikzpicture}[remember picture,overlay]
			\fill[iiiepeNavy] (current page.south west) rectangle (current page.north east);
			\node[opacity=0.13,anchor=east] at ([xshift=0.8cm]current page.east) {\includegraphics[height=9.4cm]{\iiiepedots}};
			\fill[iiiepeAqua] ([xshift=0.85cm,yshift=2.1cm]current page.south west) rectangle ([xshift=1.05cm,yshift=6.3cm]current page.south west);
			\fill[iiiepeOrange] ([xshift=1.15cm,yshift=2.1cm]current page.south west) rectangle ([xshift=1.35cm,yshift=5.35cm]current page.south west);
		\end{tikzpicture}
		\vspace{1.8cm}
		\hspace{1.25cm}
		\begin{minipage}{0.72\paperwidth}
			{\color{white}\Huge\bfseries #1}\\[0.25cm]
			{\color{white!75}\programname}
		\end{minipage}
	\end{frame}
}

\AtBeginSection[]{
	\sectionframe{\insertsectionhead}
}

\title[Plantilla IIIEPE]{Título de la presentación}
\subtitle{\programname}
\author[\studentshort]{\texorpdfstring{\studentname\\\footnotesize Matrícula: \studentid\\\href{mailto:\studentemail}{\studentemail}}{\studentname, Matricula: \studentid}}
\institute[\institutionname]{\institutionlongname\\\modality}
\date[\today]{\institutionplace\\\today}

\begin{document}

% Portada institucional
\begin{frame}[plain]
	\begin{tikzpicture}[remember picture,overlay]
		\node[anchor=center,opacity=0.18] at (current page.center) {\includegraphics[width=\paperwidth,height=\paperheight]{\iiiepecover}};
		\fill[iiiepeNavy,opacity=0.94] (current page.south west) rectangle ([xshift=0.55\paperwidth]current page.north west);
		\fill[iiiepeBlue,opacity=0.92] ([xshift=0.55\paperwidth]current page.south west) rectangle (current page.north east);
		\node[anchor=north west] at ([xshift=0.55cm,yshift=-0.45cm]current page.north west) {\includegraphics[width=4.15cm]{\iiiepelogo}};
		\node[opacity=0.17,anchor=south east] at ([xshift=0.9cm,yshift=-0.3cm]current page.south east) {\includegraphics[height=8.8cm]{\iiiepedots}};
		\fill[iiiepeOrange] ([xshift=0.56\paperwidth,yshift=1.15cm]current page.south west) rectangle ([xshift=0.92\paperwidth,yshift=1.28cm]current page.south west);
		\fill[iiiepeAqua] ([xshift=0.56\paperwidth,yshift=0.88cm]current page.south west) rectangle ([xshift=0.82\paperwidth,yshift=1.01cm]current page.south west);
	\end{tikzpicture}

	\vspace{1.9cm}
	\begin{minipage}{0.50\paperwidth}
		{\color{white}\Huge\bfseries Título de la presentación}\\[0.25cm]
		{\color{white!86}\large \programname}\\[0.35cm]
		{\color{iiiepeOrange}\rule{0.82\linewidth}{1.3pt}}\\[0.45cm]
		{\color{white}\studentname}\\
		{\color{white!80}\footnotesize Matrícula: \studentid\ · \studentemail}
	\end{minipage}
	\hfill
	\begin{minipage}{0.36\paperwidth}
		\raggedleft
		{\color{white}\Large\bfseries \institutionname}\\[0.18cm]
		{\color{white!82}\footnotesize \institutionlongname}\\[0.5cm]
		{\color{white}\small \modality}\\
		{\color{white!78}\small \periodname}\\[0.5cm]
		{\color{white!74}\footnotesize \institutionplace\ · \today}\\
		{\color{white!74}\footnotesize \href{\institutionurl}{\institutionurl}}
	\end{minipage}
\end{frame}

\begin{frame}{Contenido}
	\tableofcontents
\end{frame}

\section{Identidad del programa}

\begin{frame}{Ficha académica}
	\begin{columns}[T]
		\begin{column}{0.42\textwidth}
			\begin{center}
				\includegraphics[width=0.54\linewidth]{\iiiepemark}\\[0.35cm]
				{\Large\bfseries\color{iiiepeNavy}\institutionname}\\
				{\small\color{iiiepeMuted}Innovación para la educación}
			\end{center}
		\end{column}
		\begin{column}{0.54\textwidth}
			\begin{block}{Alumno}
				\textbf{\studentname}\\
				Matrícula: \studentid\\
				Correo: \href{mailto:\studentemail}{\studentemail}
			\end{block}
			\begin{block}{Programa}
				\programname\\
				\modality\\
				\periodname
			\end{block}
		\end{column}
	\end{columns}
	\accentbar
\end{frame}

\begin{frame}{Programa educativo cursado}
	\begin{columns}[T]
		\begin{column}{0.48\textwidth}
			\begin{block}{Datos generales}
				\begin{itemize}
					\item Programa: \programname.
					\item Sigla de referencia: \programshort.
					\item Duración: 6 tetramestres.
					\item Modalidad: presencial.
					\item Dirigido a docentes de educación básica.
				\end{itemize}
			\end{block}
		\end{column}
		\begin{column}{0.48\textwidth}
			\begin{block}{Sentido formativo}
				Formación profesionalizante orientada a problematizar la práctica,
				diseñar innovaciones sustentables y mejorar la enseñanza y el
				aprendizaje de las matemáticas con criterios de calidad y equidad.
			\end{block}
			{\scriptsize Fuente base: \href{\maemsource}{IIIEPE, ficha oficial MAEM}.}
		\end{column}
	\end{columns}
	\accentbar
\end{frame}

\begin{frame}{Objetivo formativo de la MAEM}
	\begin{block}{Propósito del programa}
		Formar docentes expertos en aprendizaje y enseñanza de las matemáticas,
		capaces de reflexionar críticamente sobre su práctica, diseñar
		innovaciones didácticas, implementarlas, evaluarlas y generar impacto en
		su comunidad profesional.
	\end{block}
	\begin{columns}[T]
		\begin{column}{0.32\textwidth}
			\begin{block}{Problematizar}
				Leer la práctica docente como objeto de análisis, no solo como rutina.
			\end{block}
		\end{column}
		\begin{column}{0.32\textwidth}
			\begin{block}{Innovar}
				Diseñar respuestas didácticas sustentables y contextualizadas.
			\end{block}
		\end{column}
		\begin{column}{0.32\textwidth}
			\begin{block}{Evaluar}
				Valorar efectos en aprendizaje, comprensión y mejora institucional.
			\end{block}
		\end{column}
	\end{columns}
	\accentbar
\end{frame}

\begin{frame}{Perfil de ingreso}
	\begin{itemize}
		\item Profesional de la educación: docente, formador, directivo o figura vinculada a la enseñanza.
		\item Experiencia o interés profesional en la enseñanza de las matemáticas en contextos educativos.
		\item Disposición para revisar y mejorar la práctica docente mediante metodologías activas e innovación.
		\item Apertura al uso de herramientas tecnológicas con sentido pedagógico.
		\item Interés por la investigación educativa aplicada y por estrategias que favorezcan el aprendizaje significativo.
	\end{itemize}
	\vfill
	{\scriptsize Síntesis elaborada a partir de la ficha oficial MAEM del IIIEPE.}
	\accentbar
\end{frame}

\begin{frame}{Perfil de egreso}
	\begin{itemize}
		\item Diseña, implementa y evalúa estrategias pedagógicas innovadoras para la comprensión matemática.
		\item Integra herramientas digitales y metodologías basadas en resolución de problemas.
		\item Adapta enfoques didácticos a contextos educativos diversos.
		\item Participa en investigación educativa sobre enseñanza y aprendizaje de las matemáticas.
		\item Lidera procesos de formación docente y mejora del aprendizaje en su comunidad educativa.
	\end{itemize}
	\vfill
	{\scriptsize Síntesis elaborada a partir de la ficha oficial MAEM del IIIEPE.}
	\accentbar
\end{frame}

\begin{frame}{Plan de estudios oficial}
	\begin{center}
		\includegraphics[height=0.74\textheight]{\maemplanimage}
	\end{center}
	\vspace{-0.15cm}
	{\scriptsize Malla curricular publicada por el IIIEPE para la \programname.}
\end{frame}

\begin{frame}{Malla curricular resumida}
	\scriptsize
	\begin{tabular}{p{0.18\textwidth}p{0.70\textwidth}}
		\toprule
		\textbf{Tetramestre} & \textbf{Unidades curriculares} \\
		\midrule
		Primero & Seminario de temas selectos de Matemáticas I; Fundamentos para la enseñanza y el aprendizaje. \\
		Segundo & Seminario de temas selectos de Matemáticas II; Aprendizaje, cambio disposicional y construcción de conocimientos complejos. \\
		Tercero & Pedagogía de las Matemáticas I; Resignificar la enseñanza de las Matemáticas. \\
		Cuarto & Fundamentos del desarrollo humano y educación; Pedagogía de las Matemáticas II. \\
		Quinto & Construcción de ambientes para el aprendizaje de las Matemáticas I; Evaluación formativa de los aprendizajes y construcción de conocimientos complejos. \\
		Sexto & Construcción de ambientes para el aprendizaje de las Matemáticas II; Propuestas innovadoras en la enseñanza de las Matemáticas. \\
		\bottomrule
	\end{tabular}
	\vfill
	{\scriptsize Los créditos se conservan en la imagen oficial; esta tabla funciona como referencia rápida editable.}
	\accentbar
\end{frame}

\section{Base para actividades}

\begin{frame}{Ejes para orientar trabajos MAEM}
	\begin{columns}[T]
		\begin{column}{0.48\textwidth}
			\begin{block}{Pregunta central}
				¿Qué problema concreto de enseñanza o aprendizaje matemático se atiende?
			\end{block}
			\begin{block}{Criterio didáctico}
				¿Qué concepto, estrategia, ambiente, evaluación o recurso justifica la intervención?
			\end{block}
		\end{column}
		\begin{column}{0.48\textwidth}
			\begin{block}{Evidencia}
				¿Qué datos, observaciones, productos o argumentos muestran comprensión?
			\end{block}
			\begin{block}{Transferencia}
				¿Cómo se puede aplicar o mejorar en la práctica docente?
			\end{block}
		\end{column}
	\end{columns}
	\accentbar
\end{frame}

\begin{frame}{Ruta sugerida para realizar una actividad}
	\begin{enumerate}
		\item Leer la consigna y convertirla en objetivo operativo.
		\item Identificar producto esperado: ensayo, cuadro, exposición, análisis, propuesta o recurso didáctico.
		\item Seleccionar conceptos clave y fuentes de apoyo.
		\item Construir el argumento: problema, fundamento, aplicación y consecuencia.
		\item Elaborar el producto visual o textual solicitado.
		\item Cerrar con reflexión profesional y conexión con la práctica docente.
	\end{enumerate}
	\accentbar
\end{frame}

\begin{frame}{Criterios de calidad para entregar}
	\begin{itemize}
		\item Claridad: cada diapositiva debe sostener una idea principal.
		\item Pertinencia: el contenido debe responder a la actividad, no solo al tema general.
		\item Evidencia: incluir fuentes, ejemplos, casos, criterios o productos verificables.
		\item Análisis: explicar relaciones, límites, implicaciones y decisiones didácticas.
		\item Identidad profesional: vincular el aprendizaje con la mejora de la práctica docente.
	\end{itemize}
	\accentbar
\end{frame}

\section{Plantilla de actividad}

\begin{frame}{Título de la actividad}
	\begin{columns}[T]
		\begin{column}{0.62\textwidth}
			\begin{itemize}
				\item Sustituir este punto por la idea principal de la actividad.
				\item Agregar evidencias, argumentos o referencias de clase.
				\item Cerrar con una implicación para la práctica docente.
			\end{itemize}
		\end{column}
		\begin{column}{0.32\textwidth}
			\begin{block}{Enfoque}
				Aprendizaje, enseñanza e innovación educativa con énfasis en matemáticas.
			\end{block}
		\end{column}
	\end{columns}
	\accentbar
\end{frame}

\begin{frame}{Objetivo de la actividad}
	\begin{block}{Redactar objetivo}
		Escribir aquí el objetivo de la actividad en una frase propia, usando un
		verbo observable: analizar, comparar, diseñar, evaluar, explicar,
		proponer, argumentar o sintetizar.
	\end{block}
	\begin{itemize}
		\item ¿Qué se va a estudiar o producir?
		\item ¿Con qué criterio teórico o didáctico se trabajará?
		\item ¿Para qué mejora, decisión o aprendizaje servirá?
	\end{itemize}
	\accentbar
\end{frame}

\begin{frame}{Análisis}
	\begin{columns}[T]
		\begin{column}{0.48\textwidth}
			\begin{block}{Concepto clave}
				Escribir aquí la definición, postura teórica o problema central.
			\end{block}
		\end{column}
		\begin{column}{0.48\textwidth}
			\begin{block}{Aplicación}
				Relacionar el concepto con la enseñanza y aprendizaje de las matemáticas.
			\end{block}
		\end{column}
	\end{columns}
	\vfill
	\begin{center}
		{\color{iiiepeMuted}\small \institutionname\ · \programshort\ · \periodname}
	\end{center}
	\accentbar
\end{frame}

\begin{frame}{Producto solicitado}
	\begin{block}{Tipo de producto}
		Indicar si la actividad pide cuadro comparativo, mapa conceptual, línea
		de tiempo, exposición, ensayo breve, análisis de caso, propuesta
		didáctica, secuencia, rúbrica o ambiente de aprendizaje.
	\end{block}
	\begin{block}{Evidencia de elaboración}
		Insertar aquí el producto construido o una síntesis visual editable.
	\end{block}
	\accentbar
\end{frame}

\begin{frame}{Aplicación a la práctica docente}
	\begin{itemize}
		\item Contexto educativo o grupo al que podría aplicarse.
		\item Problema matemático, didáctico o evaluativo que atiende.
		\item Estrategia de intervención o mejora propuesta.
		\item Evidencia esperada para valorar su impacto.
	\end{itemize}
	\accentbar
\end{frame}

\section{Cierre}

\begin{frame}{Conclusiones}
	\begin{enumerate}
		\item Sintetizar el aprendizaje principal.
		\item Explicar su relevancia para el contexto educativo.
		\item Plantear una mejora o pregunta para continuar investigando.
	\end{enumerate}
	\accentbar
\end{frame}

\begin{frame}[plain]
	\begin{tikzpicture}[remember picture,overlay]
		\fill[iiiepeNavy] (current page.south west) rectangle (current page.north east);
		\node[opacity=0.13,anchor=south east] at ([xshift=0.8cm,yshift=-0.2cm]current page.south east) {\includegraphics[height=9cm]{\iiiepedots}};
		\node[anchor=north west] at ([xshift=0.55cm,yshift=-0.45cm]current page.north west) {\includegraphics[width=3.7cm]{\iiiepelogo}};
	\end{tikzpicture}
	\vspace{2.2cm}
	\begin{center}
		{\color{white}\Huge\bfseries Gracias}\\[0.55cm]
		{\color{white}\Large \studentname}\\[0.18cm]
		{\color{white!78}\footnotesize Matrícula \studentid\ · \studentemail}\\[0.45cm]
		{\color{iiiepeOrange}\rule{0.36\paperwidth}{1.4pt}}\\[0.35cm]
		{\color{white!72}\footnotesize \institutionname\ · \programname}
	\end{center}
\end{frame}

\end{document}
